

The proof of \thmref{bayes_clt_main} will be broken up into several parts. Our
first step is to derive an explicit series expansion for posterior expectations
of the form $\expect{\post}{\phi(\theta)}$, culminating in \lemref{bayes_clt}.
In \lemref{resid_op1}, we then consider the dependence of $\phi(\theta)$ on a
subset of datapoints, writing $\phi(\theta, \xvec_s)$ (recalling
\defref{index_sets}), and showing that the results of \lemref{bayes_clt} can be
summed over $\indexset$. The residual of the expansion in \lemref{bayes_clt} is
designed so that \lemref{resid_op1} follows easily.

To prove \lemref{bayes_clt}, we follow classical proofs of the BCLT, keeping
explicit track of the residual.  We first express the desired expectation as the
ratio of integrals with respect to $\tau$.  We then divide the domain of
integration into three regions according to their distance from $\thetahat$ and
treat each separately. Finally, we combine the bounds on the integrals over
different regions for our final result.

Throughout, we will use the change of variable
%
\begin{align}
%
\theta := \thetahat + \sqrt{N} \tau
\quad\Leftrightarrow\quad
\tau = N^{-1/2} (\theta - \thetahat).
\eqlabel{tau_defn}
%
\end{align}
%
We will switch the domain of integration between $\tau$ and $\theta$ without
re-writing the arguments.  For example, for some function $\phi(\theta)$, we
will write the shorthand $\int \phi(\theta) d\tau$ for the more precise $\int
\phi\left(\sqrt{N} \tau + \thetahat\right) d\tau$. Similarly, we will overload
set notation: if $A$ is a set of $\theta$, we will write $\tau \in A$ as a
shorthand for $\tau \in \left\{\tau: \thetahat + \sqrt{N} \tau \in A \right\}$.
Although slight abuses of notation, these conventions will keep our expressions
more compact, as well as serve as a reminder of the $\theta$ domain in which the
integrand is smooth by assumption.


%%%%%%%%%%%%%%%%%%%%%%%%%%%%%%%%%%%%%%
\subsection{Ratio of integrals}\seclabel{bclt_ratio}
We need to consider the asymptotic behavior of $\expect{\post}{\phi(\theta, \xvec_s)}$,
which can be written as the ratio of the integrals
%
\begin{align*}
    %
\expect{\post}{\phi(\theta, \xvec_s)} ={}&
\frac{\int \phi(\theta, \xvec_s)
        \exp\left(N \likhat(\theta)\right)
        d \theta}
  {\int \exp\left(N\likhat(\theta)\right) d \theta}.
 %
\end{align*}
%
Recall that $\thetahat := \mathrm{argmax}_\theta \likhat(\theta)$ and
define $\thetatil(t) := t (\theta - \thetahat) + \thetahat$.  We will
employ the following Taylor series expansion:
\begin{align}
\likhat(\theta) ={}&
    \likhat(\thetahat) +
    \likhatk{1}(\thetahat)(\theta - \thetahat) +
    \frac{1}{2}\likhatk{2}(\thetahat)(\theta - \thetahat)^2 +
\nonumber
\\&
    \frac{1}{6}\likhatk{3}(\thetahat)(\theta - \thetahat)^3 +
    \tresid{4}{\likhat, \thetahat, \theta} (\theta - \thetahat)^4,
    \eqlabel{taylor_expansion}
\end{align}
%
recalling the definition of the integral form of the Taylor series residual
$\tresid{4}{\likhat, \thetahat, \theta}$ given in \defref{taylor_residual}.

Note that $\likhat(\thetahat)$ is constant in $\theta$, and that
$\likhatk{1}(\thetahat) = 0$.  Plugging in and changing the domain of
integration to $\tau$, we see
%
\begin{align*}
\expect{\post}{\phi(\theta, \xvec_s)}
={}&
\frac{\int \phi(\theta, \xvec_s)
       \exp\left(
            N \likhat(\theta) - N \likhat(\thetahat)
        \right)
       d \theta}
     {\int
     \exp\left(
        N \likhat(\theta) - N \likhat(\thetahat)
     \right) d \theta}
\\={}&
\frac{\int \phi(\theta, \xvec_s)
       \exp\left(
            N \likhat(\theta) - N \likhat(\thetahat)
        \right)
       \sqrt{N} d \tau}
     {\int
     \exp\left(
        N \likhat(\theta) - N \likhat(\thetahat)
     \right) \sqrt{N} d \tau}.
% =\\{}&
% \frac{\int \phi(\theta, \xvec_s)
%        \exp\left(
%             \frac{1}{2} \likhatk{2}(\thetahat)\tau^2 +
%             N^{-1/2} \frac{1}{6}\likhatk{3}(\thetahat) \tau^3 +
%             N^{-1} \tresid{4}{\likhat, \thetahat, \theta} \tau^4
%         \right)
%        d \tau}
%      {\int
%      \exp\left(
%          \frac{1}{2} \likhatk{2}(\thetahat)\tau^2 +
%          N^{-1/2} \frac{1}{6}\likhatk{3}(\thetahat) \tau^3 +
%          N^{-1} \tresid{4}{\likhat, \thetahat, \theta} \tau^4
%      \right) d \tau}.
\end{align*}
%
Since the factor of $\sqrt{N}$ cancels in the preceding display, it will suffice
to consider the following generic integral.
%
\begin{defn}\deflabel{bclt_integral}
%
For a target function $\phi$, define the integral functional
%
\begin{align}
I(\phi) :={}&
\int \phi(\theta, \xvec_s)
   \exp\left(
       N \likhat(\theta) - N \likhat(\thetahat)
   \right)
   d \tau \\
={}&
\int \phi(\theta, \xvec_s) \times
\\&\quad
   \exp\left(
       \frac{1}{2} \likhatk{2}(\thetahat)\tau^2 +
       N^{-1/2} \frac{1}{6}\likhatk{3}(\thetahat) \tau^3 +
       N^{-1} \tresid{4}{\likhat, \thetahat, \theta} \tau^4
   \right)
   d \tau.
   \eqlabel{integral_funcational}
\end{align}
%
\end{defn}
%
Under \defref{bclt_integral}, $\expect{\post}{\phi(\theta, \xvec_s)} = I(\phi) /
I(1)$.  We will now consider the behavior of $I(\phi)$ for a
generic third-order BCLT-okay $\phi$,
returning to consider the behavior of the ratio in \secref{bclt_region_comb}.


%%%%%%%%%%%%%%%%%%%%%%%%%%%%%%%%%%%%%%
\subsection{Regions}
For some $\delta_1$ and $\delta_2$ to be determined later, we will divide
consideration of $I(\phi)$ into the following three disjoint regions.

\begin{defn}\deflabel{region_division}
%
For any $\delta_1 > 0$ and $\delta_2 > 0$ and $N$ sufficiently large, define
the disjoint regions
%
\begin{align*}
R_1 :={}& \left\{\theta:
    \norm{\theta - \thetahat}_2 < \delta_1 \frac{\log N}{\sqrt{N}} \right\}\\
R_2 :={}& \left\{\theta:
    \delta_1 \frac{\log N}{\sqrt{N}} \le
        \norm{\theta - \thetahat}_2 < \delta_2 \right\} \\
R_3 :={}& \left\{\theta:
    \delta_2 \le \norm{\theta - \thetahat}_2  \right\}.
\end{align*}
%
Recalling our slight abuse of set notation, the same regions in the domain of $\tau$
are,
%
\begin{align*}
\tau \in R_1 \Rightarrow{}& \tau \in \left\{\tau:
    \norm{\tau}_2 < \delta_1 \log N \right\}\\
\tau \in R_2 \Rightarrow{}& \left\{\tau:
    \delta_1 \log N \le
        \norm{\tau}_2 < \delta_2 \sqrt{N} \right\} \\
\tau \in R_3 \Rightarrow{}& \left\{\tau:
    \delta_2 \sqrt{N} \le \norm{\tau}_2  \right\}.
\end{align*}
%
Similarly, we will use the regions as domains of integration for $\tau$ as well
as $\theta$. Correspondingly, we can define
%
\begin{align}
I_1(\phi) :={}& \int_{R_1} \phi(\theta, \xvec_s)
   \exp\left(N \likhat(\theta) - N \likhat(\thetahat)\right) d \tau \nonumber\\
I_2(\phi) :={}& \int_{R_2} \phi(\theta, \xvec_s)
  \exp\left(N \likhat(\theta) - N \likhat(\thetahat)\right) d \tau \nonumber\\
I_3(\phi) :={}& \int_{R_3} \phi(\theta, \xvec_s)
 \exp\left(N \likhat(\theta) - N \likhat(\thetahat)\right) d \tau.
 \eqlabel{region_integration_functional}
\end{align}
%
\end{defn}
%
Since the boundaries of the regions coincide, $I(\phi) = I_1(\phi) + I_2(\phi) +
I_3(\phi)$ for $N$ when the regions $R_1$, $R_2$, and $R_3$ are disjoint, which
will occur for $N$ sufficiently large (see \lemref{regular_event} below for more
discussion of this point). Note that the union $R_1 \bigcup R_2$ is a $\delta_2$
ball around $\thetahat$, and that $R_1$ shrinks to $\thetahat$ as $N \rightarrow
\infty$.

Throughout the proof, we will be choosing $\delta_1$ and $\delta_2$ to satisfy
certain conditions.  Our assumptions will always take the form of assuming that
$\delta_1$ is \textit{large enough}, and that $\delta_2$ is \textit{small
enough}.  For any particular $\delta_1$ and $\delta_2$, we then will require
that $N$ is large enough that the regions are disjoint.

% With the regions so defined, we will be able to choose, independently of $\xvec$,
% $\delta_1$ to be large enough and $\delta_2$ small enough that we can make the
% integral $I_1(\phi)$ contain all nonzero terms up to polynomial order in $N$
% with high probability as $N \rightarrow \infty$.  In other words, only the
% behavior of the integral $I_1(\phi)$ will matter to leading order.  Because the
% region $R_1$ is shrinking in size as $N$ grows, the integral $I_1(\phi)$ will be
% arbitrarily well approximated by a Taylor series expansion around $\thetahat$.


%%%%%%%%%%%%%%%%%%%%%%%%%%%%%%%%%%%%%%
\subsection{Nearly certain events}
Our entire analysis will be conditioned on the occurrence of several events,
each of which occurs with probability approaching one as $N$ goes to infinity.
Outside of these events we provide no guarantees.  As a consequence, the
residuals of our approximation can be used to show convergence in probability
or distribution (via Slutsky's theorem), but not convergence in expectation
nor almost sure convergence.

% \begin{defn}\deflabel{regular_event}
%

% \end{defn}

%%%%%%%%%%%%%%%%%%%%%%%%%%%%%%%%%%%%%%%%%%%
%%%%%%%%%%%%%%%%%%%%%%%%%%%%%%%%%%%%%%%%%%%

\begin{lem}\lemlabel{regular_event}
    %
For any $\epsilon_U > 0$, let $\regevent$
denote the event that all of the following occur
for a particular $N$, $\delta_1$, $\delta_2$.
(``R'' is for ``regular''.)
%
\begin{enumerate}
%
\item \eventlabel{disjoint_regions} The regions $R_1$, $R_2$, and $R_3$ are
disjoint, and $I(\phi) = I_1(\phi) + I_2(\phi) + I_3(\phi)$.
%
\item \eventlabel{thetahat_close}
$\norm{\thetahat - \thetatrue}_\infty < \epsilon_U$.
%
\item \eventlabel{prior_close}
    $|\prior(\thetahat) - \prior(\thetatrue)| < \epsilon_U$.
%
\item \eventlabel{ulln} For any scalar-valued function $\rho(\theta, \x_n)$
with $ \expect{\fdist(\xn)}{\rho(\theta, \xn)} = 0$
satisfying a ULLN with $\delta = \deltaulln$, we have
%
%\begin{align*}
%
$\sup_{\theta \in R_1 \bigcup R_2} \meann
\abs{\rho(\theta, \x_n)}
    \le \epsilon_U$.
%
\item \eventlabel{infoev} $\infoevhat > \infoev / 2$.
%
\item \eventlabel{strict_opt} There exists $\epsilon_L > 0$ such that
%
%\begin{align*}
$\sup_{\theta \in R_3}
\left( \sumn \ell(\x_n \vert \thetahat) -
       \sumn \ell(\x_n \vert \theta)
\right) > N \epsilon_L$.
%\end{align*}
\end{enumerate}

Let \assuref{bayes_clt} hold. For any $\delta_1 > 0$, any $\pi_0 \in (0,1)$, and
any $\delta_2 > 0$ and $\epsilon_U > 0$ satisfying  $\deltaulln > \epsilon_U +
\delta_2$ and $\delta_2 / 2 \ge \epsilon_U$, there exists an $N^*$ such $N >
N^*$ implies that event $\regevent$ occurs with $\fdist$-probability at least $1 -
\pi_0$.
%
\begin{proof}
%
\Eventref{disjoint_regions} follows from the fact that, for any positive
$\delta_1$ and $\delta_2$, eventually $\delta_2 \sqrt{N} > \delta_1 \log N$.

\Eventref{thetahat_close} follows from \lemref{thetahat_consistent}.

\Eventref{prior_close} follows from the fact that the prior density is
continuous by \assuref{bayes_clt} \itemref{prior_smooth} and,
by \lemref{thetahat_consistent}, we can take
$\thetahat$ to be arbitrarily close to $\thetatrue$.

For \eventref{ulln}, observe that $\norm{\thetahat - \thetatrue}_2 < \epsilon_U$
implies that
%
\begin{align*}
%
\theta \in& R_1 \bigcup R_2 \Rightarrow\\
\norm{\theta - \thetahat}_2 \le& \delta_2 \Rightarrow\\
\norm{\theta - \thetatrue}_2 \le&
    \norm{\theta - \thetahat}_2 + \norm{\thetahat - \thetatrue}_2
\le \delta_2 + \epsilon_U \le \deltaulln.
%
\end{align*}
%
Therefore,
%
\begin{align*}
%
\sup_{\theta \in R_1 \bigcup R_2} \meann \abs{\rho(\theta, \x_n)} \le&
\sup_{\theta \in \thetaball{\deltaulln}}
    \meann \abs{\rho(\theta, \x_n)} \plim 0.
%
\end{align*}
%
\Eventref{infoev} follows from \lemref{thetahat_consistent}.


For \eventref{strict_opt}, write
%
\begin{align*}
%
\MoveEqLeft
\frac{1}{N} \sup_{\theta \in R_3}
\left( \sumn \ell(\x_n \vert \thetahat) -
       \sumn \ell(\x_n \vert \theta)
\right)
=\\&
\likhat(\thetahat) - \logprior(\thetahat) -
\sup_{\theta \in R_3} \left(\likhat(\theta) - \logprior(\theta)\right)
=\\&
\likhat(\thetahat) - \likhat(\thetatrue) +
    \sup_{\theta \in R_3}
        \left(\likhat(\thetatrue) - \likhat(\theta)\right).
%
\end{align*}
%
Since $\norm{\thetahat - \thetatrue}_2 \le \epsilon_U \le \delta_2 / 2$
by assumption of the present lemma,
%
\begin{align*}
%
\theta \in& R_3 \Rightarrow\\
\norm{\theta - \thetahat + \thetahat - \thetatrue}_2 \ge & \delta_2 \Rightarrow\\
\norm{\theta - \thetahat}_2 + \norm{\thetahat - \thetatrue}_2
    \ge& \delta_2 \Rightarrow\\
\norm{\theta - \thetahat}_2 \ge& \delta_2 - \frac{\delta_2}{2} = \frac{\delta_2}{2},
%
\end{align*}
%
so the set $R_3$, is contained in $\thetadom \setminus \thetaball{\delta_2 /
2}$. Consequently, by \assuref{bayes_clt} \itemref{strict_opt}, there exists an
$\epsilon_L$ such that
%
\begin{align*}
%
\sup_{\theta: \norm{\theta - \thetahat}_2 \ge \delta_2}
    \left(\likhat(\thetatrue) - \likhat(\theta)\right) \le&
\sup_{\theta \in \thetadom \setminus \thetaball{\delta_2 / 2}}
    \left(\likhat(\thetatrue) - \likhat(\theta)\right) \le
-2 \epsilon_L.
%
\end{align*}
%
By the continuity of $\lik$, and a ULLN applied to $\likhat(\theta)$, we can
take $\norm{\thetahat - \thetatrue}_2$ sufficiently small that that
$|\likhat(\thetahat) - \likhat(\thetatrue)| < \epsilon_L$.  It follows that,
with probability approaching one,
%
\begin{align*}
%
\frac{1}{N} \sup_{\theta: \norm{\theta - \thetahat}_2 \ge \delta_2}
\left( \sumn \ell(\x_n \vert \thetahat) -
       \sumn \ell(\x_n \vert \theta)
\right) \le \epsilon_L - 2 \epsilon_L = -\epsilon_L.
%
\end{align*}
%
The result follows by multiplying both sides of the inequality by $N$.
%
\end{proof}
%
\end{lem}

Note that \lemref{regular_event} requires $\delta_2$ to be sufficiently small (so that
$\deltaulln > \delta_2$), which in turn requires $\epsilon_U$ to be small (so
that $\deltaulln > 3 \epsilon_U$).  For any particular $\deltaulln > \delta_2$,
we can always apply \lemref{regular_event} with arbitrarily small $\epsilon_U$, at the
cost of requiring a large enough $N^*$ so that the probability convergence
results apply.
%
Throughout, we will choose $\epsilon_U$ \textit{small enough} that we obtain our
desired results.



%%%%%%%%%%%%%%%%%%%%%%%%%%%%%%%%%%%%%
\subsection{Region $R_3$}

The integral over $R_3$ goes to zero by strict optimality of $\thetatrue$
and the proper behavior of the prior.

\begin{lem}\lemlabel{region_r3}
%
Under \assuref{core_bclt_assu}, when $\regevent$ occurs with
$\epsilon_U < \prior(\thetatrue) / 2$,
%
\begin{align*}
%
\norm{I_3(\phi)}_2 =
O(\exp(-N \epsilon_L)) \expect{\prior(\theta)}{\norm{ \phi(\theta, \xvec_s) }_2 }.
%
\end{align*}
%
\begin{proof}
%
We have
%
\begin{align*}
\norm{I_3(\phi)}_2
={}& \norm{ \int_{R_3} \phi(\theta, \xvec_s)
       \exp\left(
            N \likhat(\theta) - N \likhat(\thetahat)
        \right)
       d \theta }_2 \\
\le&
\sup_{\theta \in R_3}
\exp\Bigg(
    N \Bigg(
         \frac{1}{N} \sumn \ell(\x_n \vert \theta) -
         \frac{1}{N} \sumn \ell(\x_n \vert \thetahat)
     \Bigg)
 \Bigg) \times
 \int_{R_3} \norm{\phi(\theta, \xvec_s)}_2
 \frac{\prior(\theta)}{\prior(\thetahat)} d \theta \\
={}&
    \exp(-N \epsilon_L)
    \frac{1}{\prior(\thetahat)}
    \expect{\prior(\theta)}{
        \norm{\phi(\theta, \xvec_s)}_2 \mathbb{I}(\theta \in R_3)
    }
    \quad\textrm{(\lemref{regular_event}, \eventref{strict_opt})}\\
\le&
    \exp(-N \epsilon_L)
    \frac{1}{\prior(\thetatrue) - \epsilon_U}
    \expect{\prior(\theta)}{\norm{\phi(\theta, \xvec_s)}_2 \mathbb{I}(\theta \in R_3)}
\quad\textrm{(\lemref{regular_event}, \eventref{prior_close})}\\
\le&
\exp(-N \epsilon_L) 2 \prior(\thetatrue)^{-1}
\expect{\prior(\theta)}{\norm{\phi(\theta, \xvec_s)}_2 }.
\end{align*}
%
In the last step, we take $\epsilon_U < \prior(\thetatrue) / 2$ in $\regevent$
(specifically, \lemref{regular_event} \eventref{prior_close}).
%
\end{proof}
%
\end{lem}
%
We note that the form of the bound in \lemref{region_r3} is designed to
take advantage of the prior mean assumption \defref{bclt_okay} \itemref{prior_op1},
or its more general version \defref{bclt_okay_index} \itemref{prior_op1_index}.


%%%%%%%%%%%%%%%%%%%%%%%%%%%%%%%%%%%%%
\subsection{Gaussian definitions}

For regions $R_1$ and $R_2$, and the combination of the regions, we will make
use of the following compact notation for key Gaussian densities.

%%%%%%%%%%%%%%%%%%%%%%%%%%%%%%%%%%%%%%%%
%%%%%%%%%%%%%%%%%%%%%%%%%%%%%%%%%%%%%%%%
\begin{defn}\deflabel{gaussian_defs}

Write the Gaussian density with mean $0$  and covariance $\infohat^{-1}$
and its normalizing constant as
\begin{align*}
\normhatarg{\tau} :={}&
    \normalizer^{-1} \exp\left(-\frac{1}{2} \infohat \tau^2 \right)
    = \normdist(\tau \vert 0, \infohat^{-1}) \quad\textrm{ where }\quad
\normalizer :={} \det\left(2\pi \infohat^{-1} \right)^{1/2}.
\end{align*}
%
Additionally,
let $\normresid{\lambda}{S}(\tau)$ denote the $\normdist\left(\tau | 0,
\lambda^{-1} I_{\thetadim} \right)$ distribution truncated to the set $S$,
where $I_{\thetadim}$ denotes the $\thetadim \times \thetadim$ identity matrix.
Specifically,
%
\begin{align*}
%
\normresid{\lambda}{S}(\tau) :=
\frac{
    \exp\left( - \frac{\lambda}{2} \norm{\tau}_2^2 \right)
}{
    \int_{S}
        \exp\left( - \frac{\lambda}{2} \norm{\tau}_2^2 \right) d\tau
}.
%
\end{align*}
%
\end{defn}


%%%%%%%%%%%%%%%%%%%%%%%%%%%%%%%%%%%%%
\subsection{Region $R_2$}
Region $R_2$ is dealt with by choosing $\delta_2$ sufficiently small that
$I_2(\phi)$ is dominated by a particular Gaussian integral.  The we will make
that Gaussian integral small by choosing $\delta_1$ sufficiently large.

\begin{lem}\lemlabel{region_r2}
%
Let \assuref{core_bclt_assu} hold. Take $\delta_1 \ge \sqrt{8 / \infoev}$.
There exist $\delta_2$ and $\epsilon_U$, defined as functions only of the
population quantities $\infoev$ and $\likk{k}(\theta)$, such that, for any
$\xvec$ such that $\regevent$ occurs,
%
\begin{align*}
%
\norm{I_2(\phi)}_2 \le&
    O\left(N^{-3}\right)
    \expect{\normresid{\infoev / 2}{R_1 \bigcup R_2}(\tau)}
           {\norm{\phi(\theta, \xvec_s)}_2 }.
\end{align*}
%
\end{lem}
%
\begin{proof} %%%%%%%%%%%%% proof of region_r2 lemma
%
The main part of the proof consists in showing that there exists a $\delta_2$
sufficiently small so that
%
\begin{align}\eqlabel{r2_target_bound}
%
\norm{I_2(\phi)}_2 \le& \int_{R_2}
    \exp\left(
        -\frac{\infoev}{2} \norm{\tau}_2^2
        \right) \norm{\phi(\theta, \xvec_s)}_2 d\tau,
%
\end{align}
%
The approach will be to bound the difference between $N \likhat(\theta) - N
\likhat(\thetahat)$ and $\frac{1}{2} \likk{2}(\thetatrue)$ so that the entire
expression inside the integral's exponent can be written as  $A \tau^2$ for a
positive definite matrix $A$.

For the duration of this proof, define the following residual matrix:
%
\begin{align*}
\Lambda(\theta, \xvec, N) :={}&
    \frac{1}{2} \left(\likhatk{2}(\thetahat) - \likk{2}(\thetatrue)\right) + \\
&   N^{-1/2} \frac{1}{6}\left(
            \likhatk{3}(\thetahat) -
            \likk{3}(\thetatrue)
        \right) \tau + \\
&    N^{-1/2} \frac{1}{6}\likk{3}(\thetatrue) \tau +
    N^{-1} \tresid{4}{\likhat, \thetahat, \theta} \tau^2.
\end{align*}
%
Under this definition, by \eqref{taylor_expansion},
%
\begin{align*}
%
N \likhat(\theta) - N \likhat(\thetahat) ={}&
\left( \frac{1}{2} \likk{2}(\thetatrue) + \Lambda(\theta, \xvec, N) \right) \tau^2.
%
\end{align*}
%
This expansion holds for all $\theta$, but we need to control the
residual only in $R_2$.  We do so by bounding the entries of the
$\thetadim \times \thetadim$ matrix $\Lambda(\theta, \xvec,
N)$.  When $\regevent$ occurs, we have
%
\begin{align*}
%
\norm{\likhatk{2}(\thetahat) - \likk{2}(\thetatrue)}_2 \le&
\norm{\likhatk{2}(\thetahat) - \likk{2}(\thetahat)}_2 +
\norm{\likk{2}(\thetahat) - \likk{2}(\thetatrue)}_2 \\
\le&
\sup_{\theta \in R_2}
    \norm{\likhatk{2}(\theta) - \likk{2}(\theta)}_2 +
\sup_{\theta \in R_2} \norm{\likk{3}(\theta)}_2
    \norm{\thetahat - \thetatrue}_2 \\
\le& \left(1 + \sup_{\theta \in R_2} \norm{\likk{3}(\theta)}_2\right)
    \epsilon_U.
\quad\textrm{(\lemref{regular_event}, \eventref{ulln, thetahat_close})}
%
\end{align*}
%
Similarly,
%
\begin{align*}
%
\norm{\left(\likhatk{3}(\thetahat) - \likk{3}(\thetatrue)\right) \tau}_2
\le&
\norm{\likhatk{3}(\thetahat) - \likk{3}(\thetatrue)}_2
\norm{\tau}_2
\\
\le& \left(1 + \sup_{\theta \in R_2} \norm{\likk{4}(\theta)}_2\right)
    \norm{\tau}_2 \epsilon_U.
%
\end{align*}
%
Next, applying Cauchy-Schwarz gives
%
\begin{align*}
%
\norm{\likk{3}(\thetatrue) \tau}_2 \le
    \norm{\likk{3}(\thetatrue)}_2 \norm{\tau}_2.
%
\end{align*}
%
Finally, let us consider the Taylor series residual term.  Let $\thetatil(t) = t
(\theta - \thetahat) + \thetahat$.  Because we can write
%
\begin{align*}
\tresid{4}{\likhat, \thetahat, \theta} ={}&
\frac{1}{6} \int_0^1 (t - 1)^3 \likhatk{4}(\thetatil(t)) dt \\
={}&
\frac{1}{6} \int_0^1 (t - 1)^3 \left(
    \likhatk{4}(\thetatil(t)) - \likk{4}(\thetatil(t))
\right) dt +
\frac{1}{6} \int_0^1 (t - 1)^3 \likk{4}(\thetatil(t)) dt,
\end{align*}
%
we have, when $\regevent$ occurs,
%
\begin{align*}
\sup_{\theta \in R_2}
\norm{\tresid{4}{\likhat, \thetahat, \theta} \tau^2 }_2 \le&
\norm{\tresid{4}{\likhat, \thetahat, \theta}}_2 \norm{\tau}_2^2 \\
\le&
\frac{1}{6} \left(
    \sup_{\theta \in R_2}
        \norm{\likhatk{4}(\theta) - \likk{4}(\theta)}_2 +
    \sup_{\theta \in R_2} \norm{\likk{4}(\theta)}_2 \right) \norm{\tau}_2^2 \\
\le&
\frac{1}{6} \left(
    \epsilon_U + \sup_{\theta \in R_2} \norm{\likk{4}(\theta)}_2 \right)
    \norm{\tau}_2^2.
\quad\textrm{(\lemref{regular_event}, \eventref{ulln})}
\end{align*}
%

Combining the above results, together with the fact that, in $R_2$, $\sqrt{N}
\norm{\tau}_2 \le \delta_2$, gives
%
\begin{align}
%
\sup_{\theta \in R_2} \norm{\Lambda(\theta, \xvec, N)}_2 \le&
    \frac{1}{2} \norm{\likhatk{2}(\thetahat) - \likk{2}(\thetatrue)}_2 +
        \nonumber\\&
    N^{-1/2} \frac{1}{6}\norm{
        \likhatk{3}(\thetahat) -
        \likk{3}(\thetatrue)
    } \sup_{\theta \in R_2} \norm{\tau}_2 +
        \nonumber\\&
    N^{-1/2} \frac{1}{6} \norm{\likk{3}(\thetatrue)}_2
        \sup_{\theta \in R_2} \norm{\tau}_2 +
        \nonumber\\&
    N^{-1} \sup_{\theta \in R_2}
        \norm{\tresid{4}{\likhat, \thetahat, \theta}}_2 \norm{\tau}_2^2
\nonumber\\\le&
    \frac{1}{2} \left(1 + \sup_{\theta \in R_2} \norm{\likk{3}(\theta)}_2\right)
        \epsilon_U +
   \frac{1}{6} \left(1 + \sup_{\theta \in R_2} \norm{\likk{4}(\theta)}_2\right)
    \delta_2 \epsilon_U + \nonumber\\
&   \frac{1}{6} \norm{\likk{3}(\thetatrue)}_2 \delta_2 +
   \frac{1}{6} \left(
        \epsilon_U + \sup_{\theta \in R_2} \norm{\likk{4}(\theta)}_2 \right)
    \delta_2^2. \eqlabel{r2_lambda_bound}
%
\end{align}
%
The preceding bound contains only the $\epsilon_U$, $\delta_2$, and the
population quantities $\sup_{\theta \in R_2} \norm{\likk{3}(\theta)}_2$,
$\sup_{\theta \in R_2} \norm{\likk{4}(\theta)}_2$, and $\infoev$. The quantities
$\sup_{\theta \in R_2} \norm{\likk{3}(\theta)}_2$ and $\sup_{\theta \in R_2}
\norm{\likk{4}(\theta)}_2$ depend on $\delta_2$ themselves, but are finite and
decreasing in $\delta_2$.

Consequently, when $\regevent$ occurs, there exist an $\epsilon_U$ and
$\delta_2$ sufficiently small so that
%
\begin{align}\eqlabel{re2_lambda_inf_bound}
%
\norm{\Lambda(\theta, \xvec, N)}_\infty \le
    \sqrt{\thetadim^2} \norm{\Lambda(\theta, \xvec, N)}_2 \le \frac{\infoev}{2},
%
\end{align}
%
independently of $\xvec$ (except that $\regevent$ occurs). The bound
\eqref{re2_lambda_inf_bound} cannot necessarily be achieved simply by
setting $\epsilon_U$ to be small (i.e., by relying on the consistency of
$\thetahat$ and the ULLN) because of the terms $\sup_{\theta \in R_2}
\norm{\likk{4}(\theta)}_2 \delta_2^2$ and $\norm{\likk{3}(\thetatrue)}_2
\delta_2$ in \eqref{r2_lambda_bound}---it is necessary to set $\delta_2$
small as well.
% However, using the bound
% \eqref{r2_lambda_bound}, when can choose $\epsilon_U$ and $\delta_2$ as a
% function only of population quantities to guarantee \eqref{re2_lambda_inf_bound}
% whenever $\regevent$ occurs.

Recall from \defref{loglik_def} that $\info =
-\frac{1}{2}\likk{2}(\thetatrue)$, and from \defref{loglik_details_def} that
$\infoev > 0$ is the minimum eigenvalue of $\info$.  These facts, combined
with \eqref{re2_lambda_inf_bound} give that, for any $\tau$,
%
\begin{align*}
%
\sup_{\theta \in R_2}
\left(\frac{1}{2}\likk{2}(\thetatrue) + \Lambda(\theta, \xvec, N)\right) \tau^2
    \le& -\infoev \norm{\tau}_2^2 + \frac{\infoev}{2} \norm{\tau}_2^2 \\
    ={}& -\frac{\infoev}{2} \norm{\tau}_2^2.
%
\end{align*}
%
Plugging into the definition of $I_2(\phi)$ given in \defref{region_division}
gives the desired \eqref{r2_target_bound}.

To complete the proof, we use the fact that the normal density is bounded
in $R_2$.

\begin{align*}
%
\norm{I_2(\phi)}_2
\le&
\int_{R_2}
    \exp\left(
        -\frac{\infoev}{2} \norm{\tau}_2^2
        \right) \norm{\phi(\theta, \xvec_s)}_2  d\tau
\\={}&
\int_{R_2}
    \exp\left(
        -\frac{\infoev}{4} \norm{\tau}_2^2
        -\frac{\infoev}{4} \norm{\tau}_2^2
        \right) \norm{\phi(\theta, \xvec_s)}_2 d\tau
\\\le&
\sup_{\tau \in R_2}
\exp\left(
    -\frac{\infoev}{4} \norm{\tau}_2^2
\right)
\int_{R_2}
    \exp\left(
        -\frac{\infoev}{4} \norm{\tau}_2^2
        \right) \norm{\phi(\theta, \xvec_s)}_2 d\tau
\\\le &
\exp\left(
    -\frac{\infoev}{4} \delta_1^2 \log N
\right)
\int_{R_1 \bigcup R_2}
    \exp\left(
        -\frac{\infoev}{4} \norm{\tau}_2^2
        \right) \norm{\phi(\theta, \xvec_s)}_2 d\tau.
%
\end{align*}
%
In the final line, we have used the fact that $(\log N)^2 \ge \log N$ for all $N
\ge 3$. Taking $\delta_1^2 \ge 12 / \infoev$ and multiplying by the $O(1)$
normalizer for the $\normresid{\infoev/2}{R_1 \bigcup R_2}(\tau)$ completes the proof.
%
\end{proof} %%%%%%%%%%%%% proof of region_r2 lemma

By inspecting the proof of \lemref{region_r2}, one can see that that any
polynomial rate of convergence can be achieved by taking $\delta_1$ sufficiently
large. Indeed, if the lower upper boundary of $R_1$ were of the form $\tau \le
N^k$ for some $k < 1/2$, then $\norm{I_2(\phi)}_2$ would decay exponentially and
$R_1$ would still shrink to a point in the $\theta$ domain.  However, such a
rapidly growing boundary between $R_1$ and $R_2$ comes at a cost in the analysis
of region $R_1$.  In particular, it is the use of a logarithmically growing
boundary is what allows us to give polylog bounds in the subsequent
\lemref{region_r1_final_bound} and in the final result of
\thmref{bayes_clt_main}.


%%%%%%%%%%%%%%%%%%%%%%%%%%%%%%%%%%%%%
\subsection{Region $R_1$}\seclabel{region_r1}
It is region $R_1$ that will dominate the value of the integral $I(\phi)$.  We
will first perform a Taylor series expansion of all the terms around their
sample versions, and then take probability limits of the resulting expressions.

First, we will show that integrals over $\normhat$ in $R_1$ are sufficiently
close to the corresponding Gaussian integrals over $\rdom{\thetadim}$ for
sufficiently large $\delta_1$.

%%%%%%%%%%%%%%%%%%%%%%%%%%%%%%%%%%%%%%%%%%%%%%%%%%%%%%%%%%%%%%%%%%%%%%%%%
%%%%%%%%%%%%%%%%%%%%%%%%%%%%%%%%%%%%%%%%%%%%%%%%%%%%%%%%%%%%%%%%%%%%%%%%%
%%%%%%%%%%%%%%%%%%%%%%%%%%%%%%%%%%%%%%%%%%%%%%%%%%%%%%%%%%%%%%%%%%%%%%%%%

\begin{lem}\lemlabel{region_r1_gaussian_moments}
%
Let \assuref{core_bclt_assu} hold.  Let $\tau^k$ represent the
$\thetadim$-dimensional array of outer products of $\tau$.  When $\regevent$ occurs and
$\delta_1 \ge \sqrt{24 / \infoev}$ then, for any $k \ge 0$,
%
\begin{align*}
%
\int_{R_1} \normhatarg{\tau} \tau^k d\tau =
    \expect{\normhat}{\tau^k} + O(N^{-3}).
%
\end{align*}
%
\begin{proof}
%
By definition,
%
\begin{align*}
%
\int_{R_1} \normhatarg{\tau} \tau^k d\tau  - \expect{\normhat}{\tau^k}
={}&
\int_{R_2 \bigcup R_3} \normhatarg{\tau} \tau^k d\tau
%
\end{align*}
%
We will now use a technique similar to that of the proof of
\lemref{region_r2} to control the norm of the final integral.
Recall that, on $\regevent$, the minimum eigenvalue of $\infohat$, $\infoevhat$,
is at least $\frac{1}{2} \infoev$ (\lemref{regular_event} \eventref{infoev}),
so that $\normalizer^{-1} = O(1)$.
%
\begin{align*}
%
\MoveEqLeft
\norm{
    \int_{R_2 \bigcup R_3}
            \normhatarg{\tau} \tau^k d\tau
}_2
\\\le&
    \int_{R_2 \bigcup R_3}
            \normhatarg{\tau} \norm{\tau}^k_2 d\tau
\\={}&
\normalizer^{-1} \int_{R_2 \bigcup R_3}
        \exp\left(-\frac{1}{2} \infohat \tau^2 \right)
            \norm{\tau}^k_2 d\tau
\\\le&
\normalizer^{-1} \int_{R_2 \bigcup R_3}
        \exp\left(-\frac{1}{4} \infoev \norm{\tau}_2^2 \right)
            \norm{\tau}^k_2 d\tau
\\\le&
\sup_{\tau \in R_2 \bigcup R_3}
    \exp\left(-\frac{1}{8} \infoev \norm{\tau}_2^2 \right)
\normalizer^{-1} \int_{R_2 \bigcup R_3}
        \exp\left(-\frac{1}{8} \infoev \norm{\tau}_2^2 \right)
            \norm{\tau}^k_2 d\tau
\\={}&
\exp\left(-\frac{1}{8} \infoev \delta_1^2 (\log N)^2 \right)
\normalizer^{-1} \int_{R_2 \bigcup R_3}
        \exp\left(-\frac{1}{8} \infoev \norm{\tau}_2^2 \right)
            \norm{\tau}^k_2 d\tau
\\={}& O(N^{-3}),
%
\end{align*}
%
where the final line takes $\delta_1^2 \ge \frac{24}{\infoev}$, that $\log N <
(\log N)^2$ for $N \ge 3$, that $\normalizer^{-1} = O(1)$, and that normal
moments are finite.
%
\end{proof}
%
\end{lem}

%%%%%%%%%%%%%%%%%%%%%%%%%%%%%%%%%%%%%%%%%%
%%%%%%%%%%%%%%%%%%%%%%%%%%%%%%%%%%%%%%%%%%
%%%%%%%%%%%%%%%%%%%%%%%%%%%%%%%%%%%%%%%%%%

Next, we series expand the likelihood term.

% Note that you need a standalone lemma for the log likelihood expansion
% to deal with the normalizing constant.
\begin{lem}\lemlabel{region_r1_loglik_expansion}
%
Let \assuref{core_bclt_assu} hold.  For any $\xvec$ such that $\regevent$
occurs, and for $\theta \in R_1$,
%
\begin{align*}
%
\MoveEqLeft
\exp\left(N \likhat(\theta) - N \likhat(\thetahat)\right) = \\
& \exp\left(\frac{1}{2} \infohat \tau^2 \right) \times
\left( 1 + N^{-1/2} \frac{1}{6} \likhatk{3}(\theta) \tau^3 +
    \ordlog{N^{-1}} \right).
%
\end{align*}
%
\begin{proof}
%

Using \eqref{taylor_expansion}, write the log likelihood term as
%
\begin{align}
%
\MoveEqLeft
\exp\left(N \likhat(\theta) - N \likhat(\thetahat)\right) =  \nonumber \\
& \exp\left(\frac{1}{2} \likhatk{2}(\thetahat) \tau^2 \right) \times \nonumber \\
& \exp\left(
        N^{-1/2} \frac{1}{6} \likhatk{3}(\thetahat) \tau^3 +
        N^{-1} \frac{1}{6} \tresid{4}{\likhat(\theta), \thetahat, \theta} \tau^4
    \right). \eqlabel{r1_loglike_term_factorization}
%
\end{align}
%
We will expand the second term of \eqref{r1_loglike_term_factorization} using
the expansion of the  exponential function around 0, which holds for all
$z$:
%
\begin{align*}
\exp(z) ={}& 1 + z + z^2 + \frac{1}{2} \int_0^1 (1 - t)^2 \exp(tz) dt z^3 \\
={}& 1 + z + O(|z|^2).
\end{align*}
%
To match \eqref{r1_loglike_term_factorization}, we need to evaluate at
%
\begin{align*}
%
z = N^{-1/2} \frac{1}{6} \likhatk{3}(\thetahat) \tau^3 +
    N^{-1} \tresid{4}{\likhat(\theta), \thetahat, \theta} \tau^4.
%
\end{align*}
%
When $\regevent$ occurs (\lemref{regular_event} \eventref{ulln}), we
have
%
\begin{align*}
%
\norm{N^{-1} \tresid{4}{\likhat(\theta), \thetahat, \theta} \tau^4}_2 \le&
N^{-1} \sup_{\theta \in R_1}
    \norm{\tresid{4}{\likhat(\theta), \thetahat, \theta}}_2
    \sup_{\tau \in R_1}\norm{\tau^3}_2 \\
={}& \ordlog{N^{-1}}\\
%
\norm{N^{-1/2} \frac{1}{6} \likhatk{3}(\thetahat) \tau^3}_2 \le&
N^{-1/2} \frac{1}{6} \norm{\likhatk{3}(\thetahat)}_2
    \sup_{\tau \in R_1}\norm{\tau^3}_2\\
={}& \ordlog{N^{-1/2}} \Rightarrow\\
|z| ={}& \ordlog{N^{-1/2}},
%
\end{align*}
%
so that
%
\begin{align}
%
\exp\left(
        N^{-1/2} \frac{1}{6} \likhatk{3}(\thetahat) \tau^3 +
        N^{-1} \frac{1}{6} \tresid{4}{\likhat(\theta), \thetahat, \theta} \tau^4
    \right) = \nonumber\\
1 + N^{-1/2} \frac{1}{6} \likhatk{3}(\theta) \tau^3 +
    \ordlog{N^{-1}}
. \eqlabel{r1_exp_expansion}
%
\end{align}
%
Combining \eqref{r1_loglike_term_factorization, r1_exp_expansion} and
recognizing the definition of $\infohat$ gives the desired result.
%
\end{proof}
%
\end{lem}


%%%%%%%%%%%%%%%%%%%%%%%%%%%%%%%%%%%%%%%%%%%%%%%%%%%%%%%%%%%%
%%%%%%%%%%%%%%%%%%%%%%%%%%%%%%%%%%%%%%%%%%%%%%%%%%%%%%%%%%%%
%%%%%%%%%%%%%%%%%%%%%%%%%%%%%%%%%%%%%%%%%%%%%%%%%%%%%%%%%%%%

\begin{lem}\lemlabel{r1_normalizer_expansion}
%
Let \assuref{core_bclt_assu} hold.  When $\regevent$ occurs and $\delta_1$ is as
given in \lemref{region_r1_gaussian_moments},
%
\begin{align*}
%
\int_{R_1} \exp\left(
    N \likhat(\theta) - N \likhat(\thetahat)
\right) d\tau = \normalizer + \ordlog{N^{-1}}.
%
\end{align*}
%
\begin{proof}
%
Recall that, by definition,
%
\begin{align*}
%
\normalizer = \int \exp\left(-\frac{1}{2} \infohat \tau^2 \right) d\tau.
%
\end{align*}
%
Expanding using \lemref{region_r1_loglik_expansion} gives
%
\begin{align*}
%
\MoveEqLeft
\left| \int_{R_1} \exp\left(
    N \likhat(\theta) - N \likhat(\thetahat)
\right) d\tau - \normalizer \right|
\\={}&
\left|
\int_{R_1} \exp\left(\frac{1}{2} \infohat \tau^2 \right)
\left( 1 + N^{-1/2} \frac{1}{6} \likhatk{3}(\theta) \tau^3 +
    \ordlog{N^{-1}} \right) d\tau -
\int \exp\left(-\frac{1}{2} \infohat \tau^2 \right) d\tau
\right|
%
\\={}&
\left|
\int_{R_1} \exp\left(\frac{1}{2} \infohat \tau^2 \right)
\left( N^{-1/2} \frac{1}{6} \likhatk{3}(\theta) \tau^3 +
    \ordlog{N^{-1}} \right) d\tau -
\int_{R_2 \bigcup R_3} \exp\left(-\frac{1}{2} \infohat \tau^2 \right) d\tau
\right|
%
\\\le&
\frac{1}{6} N^{-1/2} \sup_{\theta \in R_1} \norm{
    \likhatk{3}(\theta) }_2
\expect{\normhat}{\ind{\tau \in R_1}\tau^3} +
\sup_{\tau \in R_2 \bigcup R_3}
    \exp\left(-\frac{1}{4} \infoev \norm{\tau}_2^2 \right) +
\ordlog{N^{-1}}
%
\\={}&
O(N^{-2}) +
\exp\left(-\frac{1}{4} \infoev \delta_1^2 (\log N)^2 \right) +
\ordlog{N^{-1}}
\\={}& \ordlog{N^{-1}},
\end{align*}
%
where, in the penultimate line, we use \lemref{region_r1_gaussian_moments},
including the given lower bound on $\delta_1$.
%
\end{proof}
%
\end{lem}

%%%%%%%%%%%%%%%%%%%%%%%%%%%%%%%%%%%%%%%%%%%%%%%%%%%%%%%%%%%%



%%%%%%%%%%%%%%%%%%%%%%%%%%%%%%%%%%%%%%%%%%
%%%%%%%%%%%%%%%%%%%%%%%%%%%%%%%%%%%%%%%%%%
%%%%%%%%%%%%%%%%%%%%%%%%%%%%%%%%%%%%%%%%%%




In \lemref{region_r1_final_bound} to follow, we integrate the result of
\lemref{region_r1_loglik_expansion} using \lemref{region_r1_gaussian_moments} to
produce our final results for region $R_1$. Recall from
\defref{loglik_details_def} that $\normhat = \normalizer^{-1}
\exp\left(\frac{1}{2} \infohat \tau^2 \right)$ is the density of the Gaussian
distribution on $\tau$ with mean zero and covariance $\infohat^{-1} =
(-\likhatk{2}(\thetahat))^{-1}$.

Note that the statement of \lemref{region_r1_final_bound} combines vector-valued
quantities like $\phigrad{2}(\thetahat, \xvec_s) \expect{\normhat}{\tau^2}$ with
scalar--valued quantities like $\norm{\phigrad{2}(\thetahat, \xvec_s)}_2$.  This is
for notational convenience since the scalar-valued terms will simply become part
of the residual, and their component values do not matter.  Formally, the
equation can be taken to hold componentwise.
%
\begin{lem}\lemlabel{region_r1_final_bound}
%
Let \assuref{core_bclt_assu} hold.  When $\delta_1$ is as given in
\lemref{region_r2}, then, for any $\xvec$ such that $\regevent$
occurs,
%
\begin{align*}
%
\MoveEqLeft
I_1(\phi) =
    \phi(\thetahat, \xvec_s)
    \int_{R_1} \exp\left(
        N \likhat(\theta) - N \likhat(\thetahat)
    \right) d\tau +
\nonumber\\ &\quad
    N^{-1} \normalizer \left(
        \frac{1}{2} \phigrad{2}(\thetahat, \xvec_s) \expect{\normhat}{\tau^2} +
        \frac{1}{6} \likhatk{3}(\thetahat) \phigrad{1}(\thetahat, \xvec_s)
            \expect{\normhat}{\tau^4} \right) +
\nonumber\\ &\quad
    \ordlog{N^{-2}} \bigg(
        1 +
        \norm{\phigrad{1}(\thetahat, \xvec_s)}_2 +
        \norm{\phigrad{2}(\thetahat, \xvec_s)}_2 +
        \expect{\normresid{\infoev / 2}{R_1}(\tau)}{
            \norm{\tresid{3}{\phi(\cdot, \xvec_s), \theta, \thetahat}}_2}
    \bigg).
%
\end{align*}
%
\begin{proof}
%
The result will follow by Taylor expanding $\phi$, combining with
\lemref{region_r1_loglik_expansion}, and then integrating using the fact that
Gaussian integrals in region $R_1$ are approximately equal to Gaussian integrals
over the whole domain, as proved in \lemref{region_r1_gaussian_moments}.


%%%%%%%%%%%%%%%%%

We first Taylor expand $\phi(\theta, \xvec_s)$ around $\thetahat$ and substitute
$\tau = \sqrt{N}(\theta - \thetahat)$.
%
\begin{align}
\phi(\theta, \xvec_s) ={}&
    \phi(\thetahat, \xvec_s) +
    N^{-1/2} \phigrad{1}(\thetahat, \xvec_s) \tau + \nonumber \\
&
    N^{-1} \frac{1}{2} \phigrad{2}(\thetahat, \xvec_s) \tau^2 +
    N^{-3/2} \tresid{3}{\phi(\cdot, \xvec_s), \theta, \thetahat} \tau^3.
    \eqlabel{r1_phi_expansion}
\end{align}


%%%%%%%%%%%%%%%%%%%%%%%%

We will now combine \eqref{r1_phi_expansion} with
\lemref{region_r1_loglik_expansion}.
%
\begin{align}
\MoveEqLeft
\exp\left(N \likhat(\theta) - N \likhat(\thetahat)\right)
\left(
    \phi(\theta, \xvec_s) - \phi(\thetahat, \xvec_s)
\right) = \nonumber \\
%%%%%%%%%%%%%%%%%%%%%%%
\MoveEqLeft
\exp\left(
    \frac{1}{2} \infohat \tau^2
\right) \times \nonumber\\
&
\left(
    N^{-1/2} \phigrad{1}(\thetahat, \xvec_s) \tau +
    N^{-1} \frac{1}{2} \phigrad{2}(\thetahat, \xvec_s) \tau^2 +
    N^{-3/2} \tresid{3}{\phi(\cdot, \xvec_s), \theta, \thetahat} \tau^3
\right)  \times \nonumber\\
&
\left(
    1 +
    N^{-1/2} \frac{1}{6} \likhatk{3}(\thetahat) \tau^3 +
    \ordlog{N^{-1}}
\right) = \nonumber\\
%%%%%%%%%%%%%%%%%%%%%%%
\MoveEqLeft
\exp\left(
    \frac{1}{2} \infohat \tau^2
\right) \times \nonumber\\
&
\Bigg(
    N^{-1/2} \phigrad{1}(\thetahat, \xvec_s) \tau +
\nonumber\\ &\quad
    N^{-1} \left(
        \frac{1}{2} \phigrad{2}(\thetahat, \xvec_s) \tau^2 +
        \frac{1}{6} \likhatk{3}(\thetahat) \phigrad{1}(\thetahat, \xvec_s) \tau^4 \right) +
\nonumber\\ &\quad
    N^{-3/2} \left(
        \tresid{3}{\phi(\cdot, \xvec_s), \theta, \thetahat} \tau^3 +
        \frac{1}{12} \likhatk{3}(\thetahat)  \phigrad{2}(\thetahat, \xvec_s) \tau^5
    \right) +
\nonumber\\ &\quad
    \ordlog{N^{-3/2}} \bigg(
        \phigrad{1}(\thetahat, \xvec_s) \tau +
% \nonumber\\ &\quad\quad\quad
        N^{-1/2}  \frac{1}{2} \phigrad{2}(\thetahat, \xvec_s) \tau^2 +
        N^{-1} \tresid{3}{\phi(\cdot, \xvec_s), \theta, \thetahat} \tau^3
    \bigg)
\Bigg) = \nonumber\\
%%%%%%%%%%%%%%%%%%%%%%%
\MoveEqLeft
\exp\left(
    \frac{1}{2} \infohat \tau^2
\right) \times \nonumber\\
&
\Bigg(
    N^{-1/2} \phigrad{1}(\thetahat, \xvec_s) \tau +
% \nonumber\\ &\quad
    N^{-1} \left(
        \frac{1}{2} \phigrad{2}(\thetahat, \xvec_s) \tau^2 +
        \frac{1}{6} \likhatk{3}(\thetahat) \phigrad{1}(\thetahat, \xvec_s) \tau^4 \right) +
\nonumber\\ &\quad
    \ordlog{N^{-3/2}} \bigg(
        \phigrad{1}(\thetahat, \xvec_s) \tau +
        \left( \ord{1} \tau^5 + \ord{N^{-1/2}}  \tau^2\right)
            \phigrad{2}(\thetahat, \xvec_s) +
\nonumber\\ &\quad\quad\quad\quad\quad\quad\quad
        \ord{1} \tresid{3}{\phi(\cdot, \xvec_s), \theta, \thetahat} \tau^3
    \bigg)
\Bigg).
\eqlabel{r1_product_expansion1}
\end{align}

In the final line of \eqref{r1_product_expansion1} we absorbed some higher-order
quantities into the lower-order $\ordlog{N^{-3/2}}$ term for convenience of
expression, and used the fact that, when $\regevent$ occurs, for $k=1,2,3$,
$\norm{\likhatk{k}(\thetahat)}_2 = \ord{1}$ when $\regevent$ occurs, since
%
\begin{align*}
%
\norm{\likhatk{k}(\thetahat) - \likk{k}(\thetahat)}_2 \le&
    \sup_{\theta \in R_1} \norm{\likhatk{k}(\theta) - \likk{k}(\theta)}_2 \le
    \epsilon_U,
%
\end{align*}
%
so
%
\begin{align*}
%
\norm{\likhatk{k}(\thetahat)}_2 \le&
    \norm{\likhatk{k}(\thetahat) - \likk{k}(\thetahat)}_2 +
    \norm{\likk{k}(\thetahat) - \likk{k}(\thetatrue)}_2 +
    \norm{\likk{k}(\thetatrue)}_2 \\
\le&
    \epsilon_U +
    \sup_{\theta \in R_1} \norm{\likk{k}(\theta) - \likk{k}(\thetatrue)}_2 + \norm{\likk{k}(\thetatrue)}_2
= \ord{1},
%
\end{align*}
%
by continuity of $\likk{k}(\theta)$.

Next, we integrate both sides of \eqref{r1_product_expansion1} using
\lemref{region_r1_gaussian_moments} to replace Gaussian integrals over $R_1$
with Gaussian moments over the whole domain, up to an order $N^{-3}$ which can
be absorbed in the $\ordlog{N^{-2}}$ error term.
%
\begin{align*}
%
I_1(\phi) ={}&
\int_{R_1} \phi(\theta, \xvec_s)
   \exp\left(N \likhat(\theta) - N \likhat(\thetahat)\right) d \tau
\\
={}&
\phi(\thetahat, \xvec_s)
\int_{R_1} \exp\left(
    N \likhat(\theta) - N \likhat(\thetahat)
\right) d\tau +
\\{}&
N^{-1} \left(
    \frac{1}{2} \phigrad{2}(\thetahat, \xvec_s)
        \expect{\normhat}{\tau^2} +
    \frac{1}{6} \likhatk{3}(\thetahat) \phigrad{1}(\thetahat, \xvec_s)
        \expect{\normhat}{\tau^4} \right) +
\nonumber\\ &
    \ordlog{N^{-2}} \bigg(
        1 +
        \norm{\phigrad{1}(\thetahat, \xvec_s)}_2 +
        \norm{\phigrad{2}(\thetahat, \xvec_s)}_2 +
\nonumber\\ &
    \hspace{5em}
        \int_{R_1}
            \exp\left(-\frac{1}{2} \infohat \tau^2 \right)
            \tresid{3}{\phi(\cdot, \xvec_s), \theta, \thetahat}
            d\tau
\bigg).
%
\end{align*}

Finally, since $\tresid{3}{\phi(\cdot, \xvec_s), \theta, \thetahat}$ depends
explicitly on $\tau$, we must control this term in a more convenient form.  We
will argue as in \lemref{region_r1_gaussian_moments}.   Recall that, when
$\regevent$ occurs, $\infoevhat \ge \frac{\infoev}{2}$, so that $\normalizer =
O(1)$ and
%
\begin{align*}
%
\MoveEqLeft
\norm{\normalizer \int_{R_1} \exp\left(-\frac{1}{2} \infohat \tau^2 \right)
    \tresid{3}{\phi(\cdot, \xvec_s), \theta, \thetahat} d\tau }_2
\\\le&
\normalizer \int_{R_1} \exp\left(-\frac{1}{2} \infohat \tau^2 \right)
    \norm{\tresid{3}{\phi(\cdot, \xvec_s), \theta, \thetahat}}_2  d\tau
\\ \le&
\normalizer \int_{R_1} \exp\left(-\frac{\infoev}{4} \norm{\tau}_2^2 \right)
    \norm{\tresid{3}{\phi(\cdot, \xvec_s), \theta, \thetahat}}_2  d\tau.
% \\ \le&
% \normalizer \int_{R_1 \bigcup R_2}
%     \exp\left(-\frac{\infoev}{4} \norm{\tau}_2^2 \right)
%         \norm{\tresid{3}{\phi(\cdot, \xvec_s), \theta, \thetahat}}_2  d\tau.
%
\end{align*}
%
The result follows by multiplying by the $O(1)$ normalizer of
the $\normresid{\infoev / 2}{R_1}$ distribution.
%
\end{proof}
%
\end{lem}


%%%%%%%%%%%%%%%%%%%%%%%%%%%%%%%%%%%%%
\subsection{Combining the regions}\seclabel{bclt_region_comb}

We now combine the results of the previous three sections to evaluate
$I(\phi)$.
The accuracy of our BCLT expansion of $\expect{\post}{\phi(\theta, \xvec_s)}$ will
depend on the following residual function.
%
%%%%%%%%%%%%%%%%%%%%%%%%%%%%%%%%%%%%%%%%%
\begin{defn}\deflabel{bclt_resid}
%
Let $\infoev > 0$ denote the minimum eigenvalue of $\info$, and let $\delta > 0$
be an arbitrarily small number. For a function $\phi(\theta, \xvec_s)$ that
is third-order BCLT-okay, define the residual function
%
\begin{align*}
%
\resid{0}(\xvec_s) :={}&
    1 +
    \norm{\phi(\thetahat, \xvec_s)}_2 +
    \norm{\phigrad{1}(\thetahat, \xvec_s)}_2 +
    \norm{\phigrad{2}(\thetahat, \xvec_s)}_2
\\
%
\resid{1}(\xvec_s) :={}&
    \expect{\normresid{\infoev / 2}{R_1}(\tau)}{
        \norm{\tresid{3}{\phi(\cdot, \xvec_s), \theta, \thetahat}}_2}
\\
%
\resid{2}(\xvec_s) :={}&
    \expect{\normresid{\infoev / 2}{R_1 \bigcup R_2}(\tau)}{\norm{\phi(\theta, \xvec_s)}_2}
\\
%
\resid{3}(\xvec_s) :={}& \expect{\prior(\theta)}{\norm{\phi(\theta, \xvec_s)}_2}
%
\\
\resid{\phi}(\xvec_s) :={}&
    \resid{0}(\xvec_s) + \resid{1}(\xvec_s) + \resid{2}(\xvec_s) + \resid{3}(\xvec_s).
%
\end{align*}
%
\end{defn}
%%%%%%%%%%%%%%%%%%%%%%%%%%%%%%%%%%%%%%%%%
%%%%%%%%%%%%%%%%%%%%%%%%%%%%%%%%%%%%%%%%%



%%%%%%%%%%%%%%%%%%%%%%%%%%%%%%%%%%%%%%%%%
%%%%%%%%%%%%%%%%%%%%%%%%%%%%%%%%%%%%%%%%%
%%%%%%%%%%%%%%%%%%%%%%%%%%%%%%%%%%%%%%%%%
%
\begin{lem}\lemlabel{bayes_clt}
%
Let \assuref{bayes_clt} hold, and assume that $\phi$ is third-order BCLT-okay.
With $\fdist$-probability approaching one,
%
\begin{align*}
%
\MoveEqLeft
\expect{\post}{\phi(\theta, \xvec_s)} - \phi(\thetahat, \xvec_s) =
\\&
N^{-1} \left(
    \frac{1}{2} \phigrad{2}(\thetahat, \xvec_s) \expect{\normhat}{\tau^2} +
    \frac{1}{6} \likhatk{3}(\thetahat) \phigrad{1}(\thetahat, \xvec_s)
        \expect{\normhat}{\tau^4} \right) +
\\&
\ordlog{N^{-2}} \resid{\phi}(\xvec_s),
%
\end{align*}
%
where $\resid{\phi}(\xvec_s)$ is given in \defref{bclt_resid}.
%
\begin{proof}
%
First, combine the results of \lemref{region_r3, region_r2, region_r1_final_bound},
consolidating higher orders into lower orders and expanding the domain of
integration in the residual $\resid{\phi}(\xvec_s)$.
%
\begin{align}
%
I(\phi) ={}&
    \phi(\thetahat, \xvec_s)
    \int_{R_1} \exp\left(
        N \likhat(\theta) - N \likhat(\thetahat)
    \right) d\tau +
\nonumber\\ &\quad
    N^{-1} \normalizer \left(
        \frac{1}{2} \phigrad{2}(\thetahat, \xvec_s) \expect{\normhat}{\tau^2} +
        \frac{1}{6} \likhatk{3}(\thetahat) \phigrad{1}(\thetahat, \xvec_s)
            \expect{\normhat}{\tau^4} \right) +
\nonumber\\ &\quad
    \ordlog{N^{-2}} \resid{\phi}(\xvec_s).
\eqlabel{i_phi_expansion}
%
\end{align}
%
For convenience in the final bound of \thmref{bayes_clt_main} we have included the
term $\norm{\phi(\thetahat, \xvec_s)}_2$ in the definition of the residual even
though it does not occur in \lemref{region_r1_final_bound}.

The final expectation is given by the ratio
%
\begin{align*}
\expect{\post}{\phi(\theta, \xvec_s)} ={}& \frac{I(\phi)}{I(1)}.
\end{align*}
%
For the remainder of this proof, for compactness of notation, we define
%
\begin{align*}
%
\zeta :={}&
    \int_{R_1} \exp\left(N \likhat(\theta) - N \likhat(\thetahat) \right) d\tau \\
={}& \normalizer + \ordlog{N^{-1}} \\
={}& \normalizer (1 + \ordlog{N^{-1}}),
%
\end{align*}
%
which follows from \lemref{r1_normalizer_expansion} and the fact that
$\normalizer$ is bounded away from zero when $\regevent$ occurs. When
$\phi(\theta, \xvec_s) \equiv 1$, then $\resid{0}(\xvec_s) = \resid{1}(\xvec_s) = 0$,
so by \lemref{region_r3, region_r2},
%
\begin{align*}
%
I(1) = \zeta\left(1 + \ord{N^{-3}}\right).
%
\end{align*}
%
We now use the expansion, valid for $|z| < 1$ as $z \rightarrow 0$,
%
\begin{align}\eqlabel{inverse_expansion}
    \frac{1}{1 + z} = 1 - z + O(z^2).
\end{align}
%
Using \eqref{i_phi_expansion} for $I(\phi)$ and
applying \eqref{inverse_expansion} to $1 / I(1)$ and to $1 / \zeta$
gives
%
\begin{align*}
%
\frac{I(\phi)}{I(1)} ={}&
\frac{I(\phi)}{\zeta( 1 + \ord{N^{-3}})} \\
={}& \frac{I(\phi)}{\zeta} + \ord{N^{-3}} I(\phi) \\
={}&
\phi(\thetahat, \xvec_s) +
N^{-1} \frac{\normalizer}{\zeta} \left(
    \frac{1}{2} \phigrad{2}(\thetahat, \xvec_s) \expect{\normhat}{\tau^2} +
    \frac{1}{6} \likhatk{3}(\thetahat) \phigrad{1}(\thetahat, \xvec_s)
        \expect{\normhat}{\tau^4} \right) +
\\&
\ordlog{N^{-2}} \resid{\phi}(\xvec_s)
\\={}&
\phi(\thetahat, \xvec_s) +
N^{-1} \left(
    \frac{1}{2} \phigrad{2}(\thetahat, \xvec_s) \expect{\normhat}{\tau^2} +
    \frac{1}{6} \likhatk{3}(\thetahat) \phigrad{1}(\thetahat, \xvec_s)
        \expect{\normhat}{\tau^4} \right) +
\\&
\ordlog{N^{-2}} \resid{\phi}(\xvec_s).
%
\end{align*}
%
Note that we have absorbed terms from the expansion of $I(\phi)$ that are linear
in $|\phi(\thetahat, \xvec_s)|$, $\norm{\phigrad{2}(\thetahat, \xvec_s)}$, and
$\norm{\phigrad{3}(\thetahat, \xvec_s)}$ into the residual $\resid{\phi}(\xvec_s)$.
To do so we have used the fact that, when $\regevent$ occurs,
$\likhatk{3}(\thetahat) = \likhatk{3}(\thetatrue) + O(1)$.

Finally, the expansion holds with $\fdist$-probability approaching one by
\lemref{regular_event}.
%
\end{proof}
%
\end{lem}
%
%%%%%%%%%%%%%%%%%%%%%%%%%%%%%%%%%%%%%%%%%


%%%%%%%%%%%%%%%%%%%%%%%%%%%%%%%%%%%%%
\subsection{Proof that the BCLT residual is square-summable}\applabel{resid_proofs}

All that remains to prove \thmref{bayes_clt_main} is to show that the residual is
$O_p(1)$ when summed.  The form of $\resid{\phi}(\xvec_s)$ is awkward, but it is
designed specifically to commute with suprema over sums, allowing us to apply
ULLNs to the remainder of our BCLT, as stated formally in
\lemref{integral_commuting_sums,taylor_commuting_sums} below.

Observe that residual expressions of the form $\sup_{\theta \in R_1 \bigcup R_2}
\norm{\phi(\theta, \x)}_2$ would not be as useful, since
%
\begin{align*}
%
\meann \sup_{\theta \in R_1 \bigcup R_2} \norm{\phi(\theta, \x_n)}_2 \ge
\sup_{\theta \in R_1 \bigcup R_2} \meann \norm{\phi(\theta, \x_n)}_2,
\end{align*}
%
and ULLN bound on the latter does not imply a bound on the former, since the
supremum may depend on $\x_n$ in a way that causes the sample mean to diverge.

%%%%%%%%%%%%%%%%%%%%%%%%%%%%%%%%%%%%%%%%%%%%%%%%%%%%%%%%%%%%%%%%%%%%%%%
%%%%%%%%%%%%%%%%%%%%%%%%%%%%%%%%%%%%%%%%%%%%%%%%%%%%%%%%%%%%%%%%%%%%%%%
%%%%%%%%%%%%%%%%%%%%%%%%%%%%%%%%%%%%%%%%%%%%%%%%%%%%%%%%%%%%%%%%%%%%%%%

\begin{lem}\lemlabel{integral_commuting_sums}
%
The Gaussian integrals of \defref{bclt_resid} commutes with sums.
%
For a scalar--valued function $\phi(\theta, \xvec_s)$ and an index set
$\indexset$,
%
\begin{align*}
%
\means \expect{\normresid{\infoev / 2}{R_1 \bigcup R_2}(\tau)}
              {\phi(\theta, \xvec_s)}^2
% \means \resid{2}(\xvec_s)
            \le
\sup_{\theta \in R_1 \bigcup R_2} \means \phi(\theta, \xvec_s)^2.
%
\end{align*}
%
\begin{proof}
%
\begin{align*}
%
\means \expect{\normresid{\infoev / 2}{R_1 \bigcup R_2}(\tau)}
              {\phi(\theta, \xvec_s)}^2
% \means \resid{2}(\xvec_s)
% ={}&
% \means \expect{\normresid{\infoev / 2}{R_1 \bigcup R_2}(\tau)}
%             {\phi(\theta, \xvec_s)^2}
\le{}&
\expect{\normresid{\infoev / 2}{R_1 \bigcup R_2}(\tau)}
              {\means \phi(\theta, \xvec_s)^2}
\\\le{}&
\sup_{\tau \in R_1 \bigcup R_2}
    \means \phi(\thetahat + \tau / \sqrt{N}, \xvec_s)^2
\\={}&
\sup_{\theta \in R_1 \bigcup R_2} \means \phi(\theta, \xvec_s)^2.
%
\end{align*}
%
\end{proof}
%
\end{lem}
%%%%%%%%%%%%%%%%%%%%%%%%%%%%%%%%%%%%%%%%%%%%%%%%%%%%%%%%%%%%%%%%%%%%%%%


%%%%%%%%%%%%%%%%%%%%%%%%%%%%%%%%%%%%%%%%%%%%%%%%%%%%%%%%%%%%%%%%%%%%%%%
%%%%%%%%%%%%%%%%%%%%%%%%%%%%%%%%%%%%%%%%%%%%%%%%%%%%%%%%%%%%%%%%%%%%%%%
%%%%%%%%%%%%%%%%%%%%%%%%%%%%%%%%%%%%%%%%%%%%%%%%%%%%%%%%%%%%%%%%%%%%%%%


\begin{lem}\lemlabel{taylor_commuting_sums}
%
The Taylor series residual of \defref{bclt_resid} commutes with sums.
%
When $\theta \in R_1 \bigcup R_2$, for a k-th order continuously differentiable
function $\phi(\theta, \xvec_s)$ and an index set $\indexset$,
%
\begin{align*}
%
\means \norm{\tresid{k}{\phi(\cdot, \xvec_s), \theta, \thetahat}}_2^2 \le
\sup_{\theta \in R_1 \bigcup R_2}
    \means \norm{\phi_{(k)}(\theta, \xvec_s)}_2^2.
%
\end{align*}
%
\begin{proof}
%
Recall from \defref{taylor_residual} that
%
\begin{align*}
%
\MoveEqLeft
\means \norm{\tresid{k}{\phi(\cdot, \xvec_s), \theta, \thetahat}}_2^2
=\\{}&
\means \norm{
    \frac{1}{(k-1)!} \int_0^1 (1 - t)^{k-1}
    \phi_{(k)}(\thetahat + t (\theta - \thetahat), \xvec_s) dt
}_2^2
\le\\&
\means \frac{1}{(k-1)!} \int_0^1
\norm{\phi_{(k)}(\thetahat + t (\theta - \thetahat), \xvec_s)}_2^2 dt
=\\&
\frac{1}{(k-1)!} \int_0^1
\means \norm{\phi_{(k)}(\thetahat + t (\theta - \thetahat), \xvec_s)}_2^2 dt
\le\\&
\frac{1}{(k-1)!}
\sup_{\theta \in R_1 \bigcup R_2}
    \means\norm{\phi_{(k)}(\theta, \xvec_s)}_2^2.
%
\end{align*}
%
The result follows because $k \ge 1$.
%
\end{proof}
%
\end{lem}

%%%%%%%%%%%%%%%%%%%%%%%%%%%%%%%%%%%%%%%%%%%%%%%%%%%%%%%%%%%%%%%%%%%%%%%


%%%%%%%%%%%%%%%%%%%%%%%%%%%%%%%%%%%%%%%%%%%%%%%%%%%%%%%%%%%%%%%%%%%%%%%
%%%%%%%%%%%%%%%%%%%%%%%%%%%%%%%%%%%%%%%%%%%%%%%%%%%%%%%%%%%%%%%%%%%%%%%
%%%%%%%%%%%%%%%%%%%%%%%%%%%%%%%%%%%%%%%%%%%%%%%%%%%%%%%%%%%%%%%%%%%%%%%

%
\begin{lem}\lemlabel{resid_op1}
%
Under \assuref{core_bclt_assu}, $\meann\resid{\phi}(\x_n)^2 = O_p(1)$.
%
\begin{proof}
%
% By \lemref{regular_event} (in particular \lemref{regular_event} \eventref{ulln}),
% together with

% TODO: it would be good to just apply the lemmas once to get a big
% expression in terms of suprema over \thetaball{\delta}.

We may take $\delta_2$ to be the smaller of the values required for
\defref{bclt_okay} \itemref{grad_op1} and
the region $R_2$ \lemref{region_r2} to
hold.  Then, by \defref{bclt_okay_index} \itemref{grad_op1_index},
for $k=0,1,2,3$,
%
\begin{align*}
%
\sup_{\theta \in R_1 \bigcup R_2}
    \meann \norm{\phigrad{k}(\theta, \x_n)}_2^2 = O_p(1),
%
\end{align*}
%

It follows immediately that $\meann
\norm{\phigrad{k}(\thetahat, \x_n)}_2 = O_p(1)$ and $\meann
\norm{\phigrad{k}(\thetahat, \x_n)}_2^2 = O_p(1)$ for $k=0, 1, 2, 3$. From this
and Cauchy-Schwarz it follows that $\meann \resid{0}(\x_n)^2 = O_p(1)$.

Applying \lemref{integral_commuting_sums} with $\indexset = \{1, \ldots, N \}$ and
$\phi(\theta, \xvec_s) = \norm{\phi(\theta, \x_n)}_2$ gives
$\meann \resid{2}(\x_n) = O_p(1)$ and $\meann \resid{2}(\x_n)^2 = O_p(1)$.

Similarly, $\meann \resid{1}(\x_n)^2 = O_p(1)$ follows from
\lemref{integral_commuting_sums} and then \lemref{taylor_commuting_sums}.

The terms $\meann \resid{3}(\x_n)$ and $\meann \resid{3}(\x_n)^2$ are controlled
by assumption in \defref{bclt_okay} \itemref{prior_op1}.

Finally, the fact that $\meann \resid{\phi}(\x_n)^2 = O_p(1)$ follows
by Cauchy-Schwarz, since $\resid{\phi}(\x_n)$ is the sum of products
of terms, the squares of each of which have $O_p(1)$ sum.
%
\end{proof}
%
\end{lem}
%
%%%%%%%%%%%%%%%%%%%%%%%%%%%%%%%%%%%%%%%%%%%%%%%%%%%%%%%%%%%%%%%%

